%% toc: P1 3 3.7,3.8
\section{Proofs for the Ultrafilter Lemma as a theorem}
\label{sec:annex-p01-w02-p01-009-ultrafilter-lemma}

Throughout, (F1)--(F4) and (D) are the conditions displayed in
\cref{sec:p01-w02-p01-009-assumptions}, and the ambient theory is
\(\mathsf{ZF}\) unless a step names \(\mathsf{AC}\).

\begin{annexitem}{annex:p01-w02-p01-009-terms}{thm:p01-w02-p01-009-terms}{Names for the two families}
\begin{proof}
\emph{The algebra \(\mathcal P(I)\).} With meet \(\cap\), join \(\cup\),
complement \(X\mapsto I\setminus X\), bottom \(\varnothing\) and top \(I\),
the set \(\mathcal P(I)\) is a Boolean algebra in the sense of
Definition~\ref{def:p01-bpi-boolean-algebra}. Each bounded distributive
lattice law and each complement law is an equality between subsets of \(I\);
it is verified by fixing \(x\in I\) and checking that the two membership
conditions for \(x\) agree, which is a classical propositional identity in
the statements \(x\in X\), \(x\in Y\), \(x\in Z\). For instance
\(x\in X\cap(Y\cup Z)\) iff \(x\in X\) and (\(x\in Y\) or \(x\in Z\)) iff
\(x\in(X\cap Y)\cup(X\cap Z)\), and \(x\in X\cap(I\setminus X)\) never holds
while \(x\in X\cup(I\setminus X)\) holds for every \(x\in I\). Its order is
inclusion: \(X\cap Y=X\) iff \(X\subseteq Y\).

\emph{Part 1.} By Definition~\ref{def:p01-bpi-filter-ideal}, a filter of
\(\mathcal P(I)\) contains the top \(I\), is upward closed under the order
\(\subseteq\) inside \(\mathcal P(I)\), and is closed under binary meets
\(\cap\); it is proper iff it omits the bottom \(\varnothing\). These four
clauses are (F1), (F4), (F3) and (F2), in that order. Hence (F1)--(F4) hold
for \(\mathcal F\) iff \(\mathcal F\) is a proper set filter on \(I\).

\emph{Part 2, from (D) to maximality.} Let \(\mathcal U\) satisfy
(F1)--(F4) and (D). By part~1 it is a proper set filter. Let \(\mathcal W\)
be a proper set filter on \(I\) with \(\mathcal U\subseteq\mathcal W\), and
let \(X\in\mathcal W\). By (D), \(X\in\mathcal U\) or
\(I\setminus X\in\mathcal U\). In the second case
\(I\setminus X\in\mathcal W\), and closure of \(\mathcal W\) under \(\cap\)
gives \(X\cap(I\setminus X)=\varnothing\in\mathcal W\), contradicting
properness of \(\mathcal W\). So \(X\in\mathcal U\). Thus
\(\mathcal W\subseteq\mathcal U\), hence \(\mathcal W=\mathcal U\), and
\(\mathcal U\) is maximal among proper set filters.

\emph{Part 2, from maximality to (D).} Let \(\mathcal U\) be a maximal
proper set filter on \(I\). By part~1 it satisfies (F1)--(F4). Fix
\(X\subseteq I\).

At most one of \(X\), \(I\setminus X\) lies in \(\mathcal U\): if both did,
(F3) would put \(X\cap(I\setminus X)=\varnothing\) in \(\mathcal U\),
against (F2).

At least one lies in \(\mathcal U\). Suppose \(X\notin\mathcal U\). By
Separation put
\[
  \mathcal W=\{\,Y\subseteq I:\ Z\cap X\subseteq Y\text{ for some }
  Z\in\mathcal U\,\}.
\]
Taking \(Z=I\) gives \(I\in\mathcal W\) and \(X\in\mathcal W\). For
\(Y\in\mathcal U\), taking \(Z=Y\) gives \(Y\cap X\subseteq Y\), so
\(\mathcal U\subseteq\mathcal W\). If \(Z\cap X\subseteq Y\subseteq Y'\subseteq I\),
then \(Z\cap X\subseteq Y'\), so \(\mathcal W\) is upward closed. If
\(Z\cap X\subseteq Y\) and \(Z'\cap X\subseteq Y'\) with
\(Z,Z'\in\mathcal U\), then \(Z\cap Z'\in\mathcal U\) by (F3) and
\((Z\cap Z')\cap X\subseteq Y\cap Y'\), so \(\mathcal W\) is closed under
\(\cap\). Thus \(\mathcal W\) is a set filter on \(I\) containing
\(\mathcal U\). It contains \(X\notin\mathcal U\), so
\(\mathcal W\ne\mathcal U\), and maximality of \(\mathcal U\) forces
\(\mathcal W\) to be improper: \(\varnothing\in\mathcal W\). So
\(Z\cap X=\varnothing\) for some \(Z\in\mathcal U\), that is,
\(Z\subseteq I\setminus X\), and (F4) gives \(I\setminus X\in\mathcal U\).
Hence (D) holds.

\emph{Part 3.} If \(I=\varnothing\), then (F1) requires
\(\varnothing=I\in\mathcal F\), while (F2) requires
\(\varnothing\notin\mathcal F\). No family meets both.

Every step instantiates finitely many given existential statements and uses
Separation; no choice principle occurs.
\end{proof}
\end{annexitem}

\begin{annexitem}{annex:p01-w02-p01-009-lemma}{thm:p01-w02-p01-009-ultrafilter-lemma}{Ultrafilter Lemma}
\begin{proof}
\emph{Part 1.} Assume \(\mathsf{ZF}+\mathsf{AC}\). Let \(I\) be a set and
let \(\mathcal F\subseteq\mathcal P(I)\) satisfy (F1)--(F4). By
Lemma~\ref{thm:p01-w02-p01-009-terms}, part~1, \(\mathcal F\) is a proper
Boolean filter of the Boolean algebra \(B=\mathcal P(I)\).
Theorem~\ref{thm:p01-bpi-ac-implication}, proved in
Annex~\ref{annex:p01-bpi-ac-implication} from \(\mathsf{AC}\) by way of
Zorn's lemma (Theorem~\ref{thm:p06-ac-to-zorn}), states that in
\(\mathsf{ZF}+\mathsf{AC}\) every proper Boolean filter is contained in an
ultrafilter, a maximal proper Boolean filter. Applied to \(B\) and
\(\mathcal F\), it gives a maximal proper set filter
\(\mathcal U\supseteq\mathcal F\) on \(I\). By
Lemma~\ref{thm:p01-w02-p01-009-terms}, part~2, \(\mathcal U\) satisfies
(F1)--(F4) and (D). This is the conclusion of part~1.

\emph{Part 2.} The derivation of part~1 is a derivation in
\(\mathsf{ZF}\) from the additional axiom \(\mathsf{AC}\), and its
conclusion holds for arbitrary \(I\) and \(\mathcal F\). Since
\(\mathsf{AC}\) is a single sentence, the deduction theorem for the
first-order ambient theory converts it into a \(\mathsf{ZF}\)-derivation of
the implication from \(\mathsf{AC}\) to the universally quantified principle
of part~1.

\emph{Part 3.} By Lemma~\ref{thm:p01-w02-p01-009-terms}, parts~1 and~2,
whose proof uses \(\mathsf{ZF}\) only, the principle of part~1 says, for every
set \(I\), that every proper set filter on \(I\) is contained in an
ultrafilter on \(I\). That is the set Ultrafilter Lemma of
Definition~\ref{def:p01-bpi-principles}, Boolean filter extension restricted
to the algebras \(\mathcal P(S)\). The two sentences are therefore
equivalent over \(\mathsf{ZF}\). Theorem~\ref{thm:p01-bpi-equivalence},
proved in \(\mathsf{ZF}\) in Annex~\ref{annex:p01-bpi-equivalence}, makes
the set Ultrafilter Lemma equivalent to \(\mathsf{BPI}\), Boolean ideal
extension and Boolean filter extension.
Theorem~\ref{thm:p01-bpi-propositional-compactness}, proved in
\(\mathsf{ZF}\) in Annex~\ref{annex:p01-bpi-propositional-compactness},
makes each of these equivalent to propositional compactness for arbitrary
sets of variables. Equivalence over \(\mathsf{ZF}\) is transitive, so all
are equivalent to the principle of part~1.

\emph{Part 4.} Assume that \(\mathsf{ZF}\) is consistent, and suppose, for
a contradiction, that \(\mathsf{ZF}\) proves the implication from the
principle of part~1 to \(\mathsf{AC}\). By part~3, \(\mathsf{ZF}\) proves the
implication from \(\mathsf{BPI}\) to that principle. Chaining the two
implications inside \(\mathsf{ZF}\), \(\mathsf{ZF}\) proves the implication
from \(\mathsf{BPI}\) to \(\mathsf{AC}\). Let
\(T=\mathsf{ZF}+\mathsf{BPI}+\neg\mathsf{AC}\). Then \(T\) proves
\(\mathsf{BPI}\) and the implication, hence \(\mathsf{AC}\) by modus ponens,
and \(T\) proves \(\neg\mathsf{AC}\) outright. So \(T\) is inconsistent.

The Halpern--L\'evy theorem, Theorem~\ref{thm:p01-bpi-choice-separation}
\cite{HalpernLevy1971}, states that if \(\mathsf{ZF}\) is consistent then
\(T\) is consistent. Since \(\mathsf{ZF}\) was assumed consistent, \(T\) is
consistent, a contradiction. Hence \(\mathsf{ZF}\) does not prove the
implication from the principle of part~1 to \(\mathsf{AC}\).

This is the only step that rests on a result whose proof is not given in this
handbook: the consistency of \(T\) relative to \(\mathsf{ZF}\) is cited, and
its model construction is not reproduced. Every other inference above is a
first-order step about provability.

\emph{Strictness.} Part~2 gives the \(\mathsf{ZF}\)-provable implication from
\(\mathsf{AC}\) to the principle, and part~4 gives, under the consistency of
\(\mathsf{ZF}\), the \(\mathsf{ZF}\)-unprovability of the converse. Their
conjunction is the stated sense in which the principle is strictly weaker
than \(\mathsf{AC}\) over \(\mathsf{ZF}\). In particular the two are not
\(\mathsf{ZF}\)-provably equivalent, so the principle is not another name for
the Axiom of Choice.
\end{proof}
\end{annexitem}

\begin{annexitem}{annex:p01-w02-p01-009-side-conditions}{thm:p01-w02-p01-009-side-conditions}{When a family of subsets can be extended}
\begin{proof}
The family \(\mathcal E^{\uparrow}\) exists by Separation from
\(\mathcal P(I)\); its defining condition quantifies over \(n\in\mathbb N\)
and finite lists from \(\mathcal E\), which are sets.

\emph{Part 1.} Let \(\mathcal U\supseteq\mathcal E\) satisfy (F1)--(F4). By
induction on \(n\in\mathbb N\) we show that for all
\(X_1,\ldots,X_n\in\mathcal E\) the intersection
\(X_1\cap\cdots\cap X_n\) belongs to \(\mathcal U\). For \(n=0\) the
intersection is \(I\), which belongs to \(\mathcal U\) by (F1). If the claim
holds for \(n\), then for \(X_1,\ldots,X_{n+1}\in\mathcal E\) we have
\(X_1\cap\cdots\cap X_n\in\mathcal U\) by the hypothesis and
\(X_{n+1}\in\mathcal E\subseteq\mathcal U\), so their intersection
\(X_1\cap\cdots\cap X_{n+1}\) belongs to \(\mathcal U\) by (F3). By (F2) no
such intersection is \(\varnothing\). So \(\mathcal E\) has the finite
intersection property. The induction uses only (F1) and (F3) for membership
and (F2) for nonemptiness.

\emph{Part 2.} Assume the finite intersection property.
\begin{itemize}
  \item \(\mathcal E\subseteq\mathcal E^{\uparrow}\): for \(X\in\mathcal E\)
    use the list of length one, \(X\subseteq X\).
  \item (F1): the empty list has intersection \(I\subseteq I\).
  \item (F2): if \(\varnothing\in\mathcal E^{\uparrow}\), some
    \(X_1\cap\cdots\cap X_n\) is a subset of \(\varnothing\), hence empty,
    against the finite intersection property.
  \item (F3): if \(X_1\cap\cdots\cap X_n\subseteq Y\) and
    \(X'_1\cap\cdots\cap X'_m\subseteq Y'\), the concatenated list
    \(X_1,\ldots,X_n,X'_1,\ldots,X'_m\) of length \(n+m\) has intersection
    \((X_1\cap\cdots\cap X_n)\cap(X'_1\cap\cdots\cap X'_m)\subseteq Y\cap Y'\),
    so \(Y\cap Y'\in\mathcal E^{\uparrow}\).
  \item (F4): if \(X_1\cap\cdots\cap X_n\subseteq Y\subseteq Y'\subseteq I\),
    then \(X_1\cap\cdots\cap X_n\subseteq Y'\).
\end{itemize}
For leastness, let \(\mathcal V\supseteq\mathcal E\) satisfy (F1)--(F4). The
induction of part~1, which used only (F1) and (F3), puts every
\(X_1\cap\cdots\cap X_n\) with \(X_i\in\mathcal E\) into \(\mathcal V\), and
(F4) then puts every \(Y\in\mathcal E^{\uparrow}\) into \(\mathcal V\). This
is the instance \(B=\mathcal P(I)\) of the generated-filter assertion of
Lemma~\ref{thm:p01-bpi-duality}, written out for set families. Each
membership check instantiates the one or two given witnessing lists; no
family of lists is chosen.

\emph{Part 3.} If \(\mathcal E\) has the finite intersection property, part~2
gives \(\mathcal E^{\uparrow}\supseteq\mathcal E\) satisfying (F1)--(F4), and
the principle of Theorem~\ref{thm:p01-w02-p01-009-ultrafilter-lemma}, part~1,
applied once with \(\mathcal F=\mathcal E^{\uparrow}\), gives
\(\mathcal U\supseteq\mathcal E^{\uparrow}\supseteq\mathcal E\) satisfying
(F1)--(F4) and (D). Conversely, if such a \(\mathcal U\) exists, part~1, which
is a \(\mathsf{ZF}\) argument, gives the finite intersection property.
\end{proof}
\end{annexitem}
